\section{The continuum form and its Green operator}\label{sec:continuum}

The electrical limit identifies the round Dirichlet integral.  We now put
this integral on the quantum area space and construct its inverse.  This
also verifies the existence and the positive-time finiteness of the
continuum heat trace that occurs in Theorem~\ref{thm:main}.

Fix a conformal sphere coordinate for the ordinary unit-area
$\sqrt 3$-quantum sphere of Proposition~\ref{prop:geometry}.  Write
$d_0$ and $\mathrm{vol}_0$ for round distance and round volume on $\Sph$,
and put
\[
 \sigma_0=\frac{\mathrm{vol}_0}{\mathrm{vol}_0(\Sph)},\qquad
 \mathcal E^0(f,g)=\int_{\Sph}\nabla f\cdot\nabla g\dd\mathrm{vol}_0,
 \qquad \mathcal E^0(f)=\mathcal E^0(f,f).
\]
We use the usual round $H^1$ norm and write $\sigma_0(f)=\int f\dd\sigma_0$.
The superscript $0$ records the absence of a factor $1/2$.  The form in
the theorem will be $\mathcal E_h=\frac12\mathcal E^0$ on its trace domain.
Our first step is a quantitative property of the area measure
$\mu_h$ which will justify that domain.

\begin{lemma}[Uniform small-ball bounds]\label{lem:quantum-ball-bounds}
Almost surely there are finite random constants $C>0$, $c>0$, and
$r_*>0$ such that, simultaneously for every $z\in\Sph$ and
$0<r\le r_*$,
\begin{equation}\label{eq:quantum-ball-bounds}
 c r^{8}\le \mu_h(B_{d_0}(z,r))\le C r^{1/56}.
\end{equation}
The constants may depend on the chosen conformal sphere coordinate.
\end{lemma}

\begin{proof}
We first prove the assertion on a fixed compact interior part of an
ordinary planar field chart.  The moment input is the subcritical GMC
theorem for a stationary positive-definite covariance
$\gamma^2\log_+(R/|x-y|)+g(x-y)$ in dimension two, with $g$ bounded and
continuous, with $0<\gamma<2$; here $\log_+u=\max\{\log u,0\}$.  For the
exact logarithmic covariance ($g=0$), its scaling
and positive and negative moment statements give
\begin{equation}\label{eq:gmc-ball-moment}
 \mathbb E[M_\gamma(B_r(z))^q]\le C_q r^{\zeta(q)},\qquad
 \zeta(q)=\left(2+\frac{\gamma^2}{2}\right)q
                    -\frac{\gamma^2}{2}q^2,
\end{equation}
uniformly for the small balls in a fixed buffer, whenever
$0<q<4/\gamma^2$ or $q<0$.  Here $M_\gamma$ is the limiting Wick chaos
with Lebesgue base measure.  This is the exact logarithmic specialization of
\cite[Propositions~3.3 and~3.5--3.7]{RobertVargas2010GMC}; the exact
logarithmic covariance is available in dimension two.  In particular the
negative moments in \eqref{eq:gmc-ball-moment} are moments of the limiting
mass of an open ball.

For $\gamma^2=3$, take $q=7/6$.  Then $q<4/3$ and $\zeta(q)=49/24$.
For $q=-1$, the exponent is $\zeta(-1)=-5$.  Markov's inequality
therefore gives, with $a=1/56$ and $b=8$,
\begin{align}
 \mathbb P\{M_\gamma(B_r(z))>r^a\}
       &\le C r^{\zeta(7/6)-(7/6)a}=C r^{97/48},
       \label{eq:gmc-upper-prob}\\
 \mathbb P\{M_\gamma(B_r(z))<r^b\}
       &\le C r^{b+\zeta(-1)}=C r^3.
       \label{eq:gmc-lower-prob}
\end{align}
At level $r=2^{-k}$, use a grid of mesh $r/16$ in the fixed buffer.
There are $O(r^{-2})$ grid points.  Apply the upper estimate to the
balls of radius $2r$ and the lower estimate to those of radius $r/4$;
constant factors in the radii only change $C$ in these displays.  The
failure probabilities summed over the grid are respectively
$O(r^{1/48})$ and $O(r)$, and hence are summable over $k$.  Almost surely,
both bounds hold on every sufficiently fine grid.  If
$r/2\le s\le r$, any ball of radius $s$ centered in the smaller compact
part is contained in a grid ball of radius $2r$ and contains a grid ball
of radius $r/4$.  Borel--Cantelli consequently gives the two polynomial
bounds for all centers and all small radii in that compact part, with
random constants.

Here is the precise passage from this reference chaos to the sphere.
A whole-plane GFF pinned to have circle average zero on a circle of
radius $R$ has covariance $\log(R/|z-w|)$ on the concentric disk of
radius $R/4$.  The field multiplied by $\gamma$ therefore has on that disk
the covariance of the exact stationary field used above.  The usual local
LQG area agrees with this
Wick chaos up to deterministic smooth normalization factors bounded above
and below on a smaller buffer.  This follows from the common-regularizer
identification, including circle averages, in
\cite[Theorem~1.1]{Berestycki2017GMC}, whose base-measure energy condition
holds for planar Lebesgue measure on a bounded buffer, and the circle-average normalization
in \cite[Propositions~1.2 and~3.2]{DuplantierSheffield2011KPZ}.
Such factors preserve the two inequalities with modified constants.
Adding any finite constant to the field also preserves them, since it
multiplies area by one positive scalar.

In the three-mark embedding of the ordinary unit-area sphere,
\cite[Proposition~3.5]{C} gives local mutual absolute continuity, modulo
additive constants, with a whole-plane GFF on bounded buffers avoiding
the three marks.  Existence of the local polynomial bounds with some
constants is invariant under adding a field constant; the grid formulation
makes it a measurable event of the local field modulo constants.  It is
therefore an almost-sure local property transferred by that proposition.  We use
it on a countable collection of nested compact and open buffers exhausting
the complement of the marks.  Resample a second triple of area marks and
use the same proposition in its three-mark coordinate.  The sphere area
is nonatomic, as in the sphere input of Proposition~\ref{prop:geometry},
so the two triples are disjoint almost surely.  Their two mark-free open
sets cover the sphere.  On each realization, select finitely many of the
buffered interior charts whose interiors cover it.  The coordinate
changes are smooth sphere conformal maps, with upper and lower derivative
bounds on these compact patches.  A Lebesgue number for the finite cover
and comparison of round and chart balls now give
\eqref{eq:quantum-ball-bounds}, with common random constants and a common
small radius.  Thus the estimate for the conditioned sphere follows from
these local Gaussian comparisons.
\end{proof}

The remaining construction is deterministic.  Its input will only be
uniform upper and lower polynomial ball bounds.  The upper bound controls
logarithmic potentials; the lower bound will make the Sobolev trace
injective.  The quasi-continuous trace description for a planar massive
field appears in \cite[Section~2.3]{AndresKajino2016Heat}, and the use of a
Hilbert--Schmidt Green operator for a log-correlated Gaussian field on a
flat torus appears in \cite[Section~3]{MaillardRhodesVargasZeitouni2016Heat}.
We give the corresponding deterministic argument for the measure at hand.

Let $G$ be the symmetric round Green kernel with zero round mean.  Its
normalization, including its scalar, is specified by
\begin{equation}\label{eq:round-green-normalization}
 \int G(x,y)\dd\sigma_0(x)=0,\qquad
 \mathcal E^0(G(\cdot,y),\phi)=\phi(y)-\sigma_0(\phi)
 \quad(\phi\in C^\infty(\Sph)).
\end{equation}
The second equality is distributional.  In a local conformal coordinate,
\[
 G(x,y)=-(2\pi)^{-1}\log|x-y|+O(1)
\]
near the diagonal.  This is the Green kernel of the no-half integral
$\mathcal E^0$.

\begin{lemma}[Sobolev trace and Green potentials]\label{lem:trace-green}
Let $\mu$ be a Borel probability measure on the round sphere.  Suppose
that, for constants $a>0$, $b<\infty$, $c,C>0$, and $r_*>0$,
\begin{equation}\label{eq:deterministic-ball-bounds}
 c r^b\le\mu(B_{d_0}(x,r))\le C r^a
 \quad(x\in\Sph, 0<r\le r_*).
\end{equation}
Smooth restriction extends uniquely to a bounded injective map
$R_\mu:H^1(\Sph,\mathrm{vol}_0)\longrightarrow L^2(\mu)$.
The measure $\mu$ charges no $H^1$-polar set, and $R_\mu f$ is the
quasi-continuous representative of $f$, $\mu$-almost everywhere.  There
is a constant $C_\mu<\infty$ such that
\begin{align}
 \|R_\mu(f-\sigma_0(f))\|_{L^2(\mu)}
       &\le C_\mu\mathcal E^0(f)^{1/2},\label{eq:centered-trace}\\
 C_\mu^{-1}\|f\|_{H^1}
       &\le \bigl(\|R_\mu f\|_{L^2(\mu)}^2+\mathcal E^0(f)\bigr)^{1/2}
       \le C_\mu\|f\|_{H^1}.\label{eq:trace-norm-equivalence}
\end{align}
Moreover, for each $g\in L^2(\mu)$ there is a unique $u_g\in H^1$ with
$\sigma_0(u_g)=0$ such that
\begin{equation}\label{eq:potential-weak-equation}
 \mathcal E^0(u_g,f)=\int g R_\mu f\dd\mu
               -\left(\int g\dd\mu\right)\sigma_0(f)
 \quad(f\in H^1).
\end{equation}
The integral operator $T_\mu$ with kernel $G$ is a Hilbert--Schmidt
operator on $L^2(\mu)$, and the potential has the trace
\begin{equation}\label{eq:potential-trace}
 R_\mu u_g=T_\mu g,\qquad
 T_\mu g(x)=\int G(x,y)g(y)\dd\mu(y)
 \quad\hbox{in }L^2(\mu).
\end{equation}
\end{lemma}

\begin{proof}
Decrease $r_*$ to at most one and set
$L(x,y)=1+\log_+(r_*/d_0(x,y))$.  The upper ball bound and layer cake
give, for $j=1,2$,
\begin{equation}\label{eq:log-potential-bound}
 \sup_x\int L(x,y)^j\dd\mu(y)<\infty.
\end{equation}
Indeed the nonconstant part is bounded using
\[
 \int\log_+\!\left(\frac{r_*}{d_0(x,y)}\right)^j\dd\mu(y)
 =j\int_0^\infty t^{j-1}\mu(B_{d_0}(x,r_*e^{-t}))\dd t
 \le Cj r_*^a\int_0^\infty t^{j-1}e^{-at}\dd t.
\]
The same bounds hold for $\sigma_0$.  Also $\mu$ is nonatomic by the
upper bound.  In particular the diagonal has zero mass for all products
of measures absolutely continuous with respect to $\mu+\sigma_0$.

For $g\in L^2(\mu)$, write $\mu(g)=\int g\dd\mu$ and consider the
zero-mass signed measure
\[
 s_g=g\mu-\mu(g)\sigma_0.
\]
The bound $|G(x,y)|\le C L(x,y)$ and Schur's inequality imply
\begin{equation}\label{eq:signed-log-energy}
 \iint L(x,y)\dd|s_g|(x)\dd|s_k|(y)
       \le C\|g\|_{L^2(\mu)}\|k\|_{L^2(\mu)}.
\end{equation}
For example the $\mu\otimes\mu$ term has this bound by the $j=1$ case
of \eqref{eq:log-potential-bound}; the mixed terms follow by integrating
$L$ against $\sigma_0$ first and using $|\mu(g)|\le\|g\|_2$.  Thus all
signed Green integrals below are absolutely defined.

We record why these integrals are Dirichlet energies, including for the
signed measures $s_g$.  Let $\Delta_0$ be the round Laplace--Beltrami
operator with nonpositive spectrum.  Thus $-\Delta_0$ is the operator of
$\mathcal E^0$ in $L^2(\mathrm{vol}_0)$.  Let $P_t=e^{t\Delta_0}$ and
let $p_t$ be its kernel relative to $\mathrm{vol}_0$.  For $t>0$, put
\[
 G_t(x,y)=\int_t^\infty
       \bigl(p_s(x,y)-\mathrm{vol}_0(\Sph)^{-1}\bigr)\dd s.
\]
The small-time Gaussian estimate for the smooth round heat kernel and
its exponential convergence to the constant kernel for large times give
\begin{equation}\label{eq:smoothed-green-bound}
 |G_t(x,y)|\le C L(x,y)\quad(t>0, x\ne y),\qquad
 G_t(x,y)\longrightarrow G(x,y)\quad(t\downarrow0, x\ne y).
\end{equation}
For instance, the singular part is controlled by
\[
 \int_0^1 s^{-1}\exp(-d_0(x,y)^2/(Cs))\dd s\le C L(x,y).
\]
Let $u_{g,t}$ be the unique smooth function with $\sigma_0(u_{g,t})=0$
and
\[
 \mathcal E^0(u_{g,t},\phi)=\int P_t\phi\dd s_g
 \quad(\phi\in C^\infty(\Sph));
\]
its source is the smoothed measure $P_t^*s_g$.  The heat semigroup identity and integration by parts give
the exact cross-energy formula
\begin{equation}\label{eq:smoothed-cross-energy}
 \mathcal E^0(u_{g,t},u_{k,s})
       =\iint G_{t+s}(x,y)\dd s_g(x)\dd s_k(y).
\end{equation}
There is no diagonal mass, so \eqref{eq:signed-log-energy} and
\eqref{eq:smoothed-green-bound} justify dominated convergence in this
formula.  Applying it to the three terms in
$\mathcal E^0(u_{g,t}-u_{g,s})$ shows that these potentials are Cauchy
in Dirichlet norm.  Poincar\'e's inequality and their zero round means
make them Cauchy in $H^1$.  Their limit $u_g$ satisfies
\begin{align}
 \sigma_0(u_g)&=0,\qquad \mathcal E^0(u_g)\le C\|g\|_2^2,
       \label{eq:potential-energy}\\
 \mathcal E^0(u_g,u_k)&=\iint G(x,y)\dd s_g(x)\dd s_k(y),
       \label{eq:cross-green-energy}\\
 \mathcal E^0(u_g,\phi)&=\int \phi\dd s_g
       \quad(\phi\in C^\infty(\Sph)).\label{eq:smooth-potential-equation}
\end{align}
The last identity follows by letting $t\downarrow0$ in the weak equation
for $u_{g,t}$, since $P_t\phi\to\phi$ uniformly.  This also identifies
the limit as a distributional potential.  Poincar\'e gives its uniqueness
among zero-round-mean $H^1$ functions.

For smooth $f$, \eqref{eq:smooth-potential-equation} and
\eqref{eq:potential-energy} show that
\[
 \left|\int (f-\sigma_0(f))g\dd\mu\right|
 \le C\|g\|_{L^2(\mu)}\mathcal E^0(f)^{1/2}.
\]
Taking the $L^2(\mu)$ dual supremum proves \eqref{eq:centered-trace} on
smooth functions.  Together with the bound on the round mean it gives
$\|f\|_{L^2(\mu)}\le C\|f\|_{H^1}$.  The density of smooth functions in
$H^1$ gives a unique bounded extension $R_\mu$ and proves
\eqref{eq:centered-trace} for every $f\in H^1$.  It also extends
\eqref{eq:smooth-potential-equation} to \eqref{eq:potential-weak-equation}.

We next identify this extension with the usual quasi-continuous values.
For a compact set $F$, the $H^1$ capacity can be defined as the infimum
of $\|\phi\|_{H^1}^2$ over smooth $\phi\ge1$ on a neighborhood of $F$,
and then extended by the usual outer regularization.  For each such
test function the restriction inequality gives
\[
 \mu(F)\le\int \phi^2\dd\mu\le C\|\phi\|_{H^1}^2.
\]
Infimizing, and then using inner regularity of $\mu$, proves that every
Borel set of zero $H^1$ capacity is $\mu$-null.  If smooth functions
$f_n$ converge to $f$ in $H^1$, the capacity Chebyshev inequality applied
to a sufficiently rapidly converging subsequence shows that this
subsequence converges outside a set of capacity zero to the
quasi-continuous representative $\widetilde f$.  A further subsequence
converges $\mu$-almost everywhere to $R_\mu f$, by its $L^2(\mu)$
convergence.  The two limits agree because the exceptional polar set is
$\mu$-null.  This proves the asserted identification of the trace.

To see why the trace is injective, work in a fixed smooth planar
coordinate patch and take a Euclidean disk $B_{3r}(x)$ contained in that
patch.  Round and Euclidean balls are uniformly comparable on its fixed
buffer, so \eqref{eq:deterministic-ball-bounds} holds there for Euclidean
balls after changing constants.  Rescale $B_r(x)$ onto $B_1$ and let
$\widehat\mu_{x,r}$ be the pushforward of
$\mu|_{B_r(x)}/\mu(B_r(x))$.  Let $\mathcal L_{B_1}$ be normalized
Lebesgue measure on $B_1$.  For $0<s\le1$ and any center $y$,
\[
 \widehat\mu_{x,r}(B_s(y))
 \le\min\{1,A_r s^a\},\qquad
 A_r=\max\{1,C_1 r^{a-b}\}.
\]
Another layer-cake calculation gives
\begin{align}
 \sup_y\int\log_+\frac1{|y-z|}\dd\widehat\mu_{x,r}(z)
 &\le\int_0^\infty\min\{1,A_r e^{-at}\}\dd t\notag\\
 &\le C(1+|\log r|).\label{eq:rescaled-log-bound}
\end{align}
The Dirichlet Green kernel of $B_3$ is nonnegative and, on
$B_1\times B_1$, is bounded by $C(1+\log_+(1/|y-z|))$.  Its analogous
bound for $\mathcal L_{B_1}$ is uniform.  The Dirichlet version of the
regularization argument \eqref{eq:smoothed-cross-energy}, applied to
$\widehat\mu_{x,r}-\mathcal L_{B_1}$, therefore yields
\begin{equation}\label{eq:local-dual-bound}
 \|\widehat\mu_{x,r}-\mathcal L_{B_1}\|_{H^{-1}(B_3;\,\|\nabla\cdot\|_2)}
       \le C\sqrt{1+|\log r|}.
\end{equation}
Here $H^{-1}(B_3;\|\nabla\cdot\|_2)$ is the dual of $H^1_0(B_3)$ with
its Dirichlet norm.  One obtains the identity by smoothing inside $B_3$
and then taking limits, just as above; the supports in $B_1$ leave a
fixed boundary buffer, and \eqref{eq:rescaled-log-bound} provides the
integrable dominating logarithm.

For smooth $f$, set $F(y)=f(x+ry)$ and choose a smooth cutoff $\chi$
supported in $B_3$ and equal to one on $B_1$.  Test
\eqref{eq:local-dual-bound} on
$\chi(F-|B_3|^{-1}\int_{B_3}F\dd y)$.  Poincar\'e's inequality on $B_3$ bounds
the Dirichlet norm of this cutoff by $C\|\nabla F\|_{L^2(B_3)}$.
The signed functional has mass zero and is supported where $\chi=1$.
Scaling back, using the two-dimensional invariance of the Dirichlet
integral, gives
\begin{equation}\label{eq:local-average-trace}
 \left|\frac{1}{|B_r|}\int_{B_r(x)} f\dd y
       -\frac{1}{\mu(B_r(x))}\int_{B_r(x)}R_\mu f\dd\mu\right|
 \le C\sqrt{1+|\log r|}
       \left(\int_{B_{3r}(x)}|\nabla f|^2\dd y\right)^{1/2}.
\end{equation}
This inequality extends to $H^1$ functions.  Indeed smooth approximation
of the cutoff converges in $H^1_0(B_3)$, and both the global restriction
map and the functional in \eqref{eq:local-dual-bound} are continuous;
their integrals agree on the smooth approximants.  If $R_\mu f=0$, the
second average in \eqref{eq:local-average-trace} is zero.  At almost
every Lebesgue point of both $f$ and $|\nabla f|^2$, its first average
tends to $f(x)$ while its right side is
$O_x(r\sqrt{1+|\log r|})$, hence tends to zero.  It follows that $f=0$
almost everywhere in each patch, and therefore on the sphere.  Thus
$R_\mu$ is injective.

For the norm comparison, write $f=c+v$ with $c=\sigma_0(f)$ and
$\sigma_0(v)=0$.  Poincar\'e and \eqref{eq:centered-trace} give
\[
 \|v\|_{H^1}\le C\mathcal E^0(f)^{1/2},\qquad
 |c|\le\|R_\mu f\|_2+\|R_\mu v\|_2
       \le\|R_\mu f\|_2+C\mathcal E^0(f)^{1/2}.
\]
These imply the left inequality in \eqref{eq:trace-norm-equivalence};
boundedness of restriction gives the right one.

It remains to verify the claimed trace of the potential.  The $j=2$
case of \eqref{eq:log-potential-bound} makes
$G\in L^2(\mu\otimes\mu)$, so $T_\mu$ in \eqref{eq:potential-trace} is
well defined and Hilbert--Schmidt.  For $g,k\in L^2(\mu)$, use
\eqref{eq:potential-weak-equation} with $u_k$ and $f=u_g$.  Since
$\sigma_0(u_g)=0$, the cross-energy identity gives
\[
 \int k R_\mu u_g\dd\mu
 =\mathcal E^0(u_k,u_g)
 =\iint G(x,y)k(x)g(y)\dd\mu(x)\dd\mu(y).
\]
In the last equality the terms involving $\sigma_0$ in
\eqref{eq:cross-green-energy} vanish by the centering
\eqref{eq:round-green-normalization}.  The displayed identity for every
$k$ proves \eqref{eq:potential-trace} in $L^2(\mu)$.
\end{proof}

\begin{proposition}[Continuum form and inverse]\label{prop:continuum-form}
Almost surely, the closure in $L^2(\mu_h)$ of the smooth form
$\frac12\mathcal E^0$ has domain
\begin{equation}\label{eq:continuum-form-domain}
 \mathcal D_h=R_{\mu_h}(H^1(\Sph,\mathrm{vol}_0)),\qquad
 \mathcal E_h(R_{\mu_h}f,R_{\mu_h}g)=\frac12\mathcal E^0(f,g).
\end{equation}
It is a densely defined closed regular Dirichlet form.  Let $A_h$ be its
nonnegative self-adjoint operator, and set $B_h=2A_h$, the operator of the
raw form $\mathcal E^0$.  Put $\mu=\mu_h$ and define
\begin{align}
 G_\mu(x,y)={}&G(x,y)-\int G(x,z)\dd\mu(z)-\int G(z,y)\dd\mu(z)
       +\iint G(z,w)\dd\mu(z)\dd\mu(w),\label{eq:double-centered-green}\\
 K_hg(x)={}&\int G_\mu(x,y)g(y)\dd\mu(y).
       \label{eq:continuum-inverse}
\end{align}
The kernel $G_\mu$ is in $L^2(\mu\otimes\mu)$ and has zero $\mu$-mean
in each variable.  The operator $K_h$ is zero on constants and is the
positive, injective inverse of $B_h$ on $L^2_0(\mu)$, the subspace of
$\mu$-mean-zero functions.

Both $B_h$ and $A_h$ have compact resolvent.  On the infinite-dimensional
space $L^2(\mu)$ their eigenvalues, counted with multiplicity, can be
written
\[
 0=\nu_0<\nu_1\le\nu_2\le\cdots\longrightarrow\infty,
 \qquad \Lambda_j=\frac{\nu_j}{2}\quad(j\ge0),
\]
respectively.  The constant eigenfunction is the only zero mode, and
\begin{equation}\label{eq:continuum-inverse-summability}
 \sum_{j\ge1}\nu_j^{-2}=\|K_h\|_{\HS}^2<\infty,
 \qquad
 \Tr e^{-tA_h}=1+\sum_{j\ge1}e^{-t\nu_j/2}<\infty\quad(t>0).
\end{equation}
The heat series converges locally uniformly on $(0,\infty)$.  The ordered
eigenvalues and this $C_{\mathrm{loc}}((0,\infty))$ heat trace have Borel
versions as functions of the area measure in the fixed sphere coordinate.
The form and these readouts are invariant under conformal or
anticonformal changes of that coordinate.
\end{proposition}

\begin{proof}
Apply Lemma~\ref{lem:trace-green} to $\mu_h$ using
Lemma~\ref{lem:quantum-ball-bounds}.  Injectivity makes the formula
\eqref{eq:continuum-form-domain} well defined.  The norm equivalence
\eqref{eq:trace-norm-equivalence} proves that this domain is complete in
the form norm.  It is exactly the closure of the smooth domain, since
smooth functions are dense in $H^1$.  Smooth restrictions are also dense
in $L^2(\mu)$: continuous functions are dense for a finite Radon measure
on a compact space, and smooth functions are uniformly dense in the
continuous ones.  Sobolev normal contractions commute with the
quasi-continuous restriction and decrease the Dirichlet integral.  This
proves the Markov property.  The same smooth functions form a uniform
core in $C(\Sph)$ and a form core, proving regularity.  The representation
theorem for closed forms gives $A_h$ and $B_h$; multiplying a form by
two multiplies its associated operator by two, so $B_h=2A_h$.

This also identifies the form with the Liouville time change of standard
conformal Brownian motion.  Take round Brownian motion with generator
$\frac12\Delta_0$ and base form $\frac12\mathcal E^0$ in
$L^2(\mathrm{vol}_0)$.  The finite measure $\mu$ charges no polar set,
so it is a smooth measure for this form.  The extended finite-energy space
of the round form is $H^1$ on this compact sphere.  Indeed, for an energy-Cauchy sequence,
subtracting round means makes it $H^1$-Cauchy by Poincar\'e; a finite
almost-everywhere limit then forces a subsequence of the subtracted
constants to converge as well.  Injectivity of the quasi-continuous trace
now gives full quasi-support: a function in this extended space whose
trace vanishes $\mu$-almost everywhere is zero in $H^1$ and hence
quasi-everywhere.  The time-change theorem
\cite[Theorem~6.2.1]{FukushimaOshimaTakeda2011} therefore gives exactly
\eqref{eq:continuum-form-domain} for the time change whose Revuz measure
is $\mu$.  In particular its Brownian clock is the half-form clock.

For later use the raw form is coercive on the centered subspace.  If
$v=R_\mu f$ and $\mu(v)=0$, then
$\sigma_0(f)=-\mu(R_\mu(f-\sigma_0(f)))$.  Hence
\begin{equation}\label{eq:continuum-centered-poincare}
 \|v\|_{L^2(\mu)}\le 2C_\mu\mathcal E^0(f)^{1/2}.
\end{equation}
The zero-energy functions are precisely the constants, by the round
Poincar\'e inequality and trace identification.

Let $\Pi_\mu g=\mu(g)\mathbf1$ and $P_\mu=\Id-\Pi_\mu$.  Double
centering in \eqref{eq:double-centered-green} is the orthogonal projection
$P_\mu\otimes P_\mu$ of $G$ in $L^2(\mu\otimes\mu)$.  Thus its two
marginals vanish and it is square integrable.  Its operator is
$K_h=P_\mu T_\mu P_\mu$, which is Hilbert--Schmidt and zero on constants.
For centered $g$, Lemma~\ref{lem:trace-green} gives
\[
 K_hg=R_\mu\bigl(u_g-\mu(R_\mu u_g)\bigr),\qquad
 \mathcal E^0(u_g-\mu(R_\mu u_g),f)=\int g R_\mu f\dd\mu
 \quad(f\in H^1).
\]
This is exactly the domain criterion for the operator of the raw closed
form: $K_hg\in\operatorname{Dom}(B_h)$ and $B_hK_hg=g$.  Uniqueness on
the centered subspace follows from
\eqref{eq:continuum-centered-poincare}.  Consequently $K_h=B_h^{-1}$
there.  In particular it is positive and injective on that subspace.
The output projection in $P_\mu T_\mu P_\mu$ chooses the representative
of the weak solution with $\mu$-mean zero.

The inverse identity gives
\[
 (B_h+\Id)^{-1}=\Pi_\mu+K_h(\Id+K_h)^{-1}.
\]
The right side is compact, since $\Pi_\mu$ has rank one and $K_h$ is
compact.  Scaling also gives compact resolvent for $A_h$.  The ball
bounds give full support and no atoms, so $L^2(\mu)$ is
infinite-dimensional.  The compact-resolvent spectral theorem and the
one-dimensional kernel now give the stated eigenvalue enumeration.
On the centered space the inverse eigenvalues are $\nu_j^{-1}$.
The Hilbert--Schmidt identity yields the first assertion in
\eqref{eq:continuum-inverse-summability}.  For every $\delta>0$,
\[
 e^{-t s/2}\le C_\delta s^{-2}
 \quad(s>0, t\ge\delta),
 \qquad C_\delta=\sup_{s>0}s^2e^{-\delta s/2}<\infty.
\]
The Weierstrass test therefore proves the heat assertion and its local
uniform convergence, with the constant mode counted once.

For completeness, choose a countable rational linear space of smooth
functions dense in $H^1$, for example rational combinations of round
spherical harmonics.  In the min--max formula for $\nu_j$ one may take
the infimum only over $(j+1)$-dimensional spans from this space: the
form-norm density just proved approximates any finite-dimensional test
space together with its mass and energy matrices.  For each fixed tuple
of smooth functions, the entries of the mass matrix are
$\int\phi_k\phi_l\dd\mu$, which are Borel, indeed continuous, in the
weak topology on probability measures; the energy matrix is deterministic.
The supremum of its Rayleigh quotient can be taken over nonzero rational
vectors, assigning the value $+\infty$ when the mass denominator vanishes.
This is a countable supremum of Borel functions.  The countable infimum
over tuples is therefore Borel and equals $\nu_j$ on the measures under
consideration.  Finite heat sums are Borel maps into
$C_{\mathrm{loc}}((0,\infty))$, and their locally uniform limit is
Borel as well.

Finally, a conformal or anticonformal sphere map transports $\mu$ and
acts unitarily on the corresponding $L^2$ spaces.  It preserves the
two-dimensional Dirichlet integral.  Being a smooth diffeomorphism of
the compact sphere, it preserves $H^1$ and capacity-zero sets, so it
also transports the quasi-continuous trace domain.  The resulting closed
forms are unitarily equivalent.  This proves the asserted coordinate
invariance and completes the continuum construction.
\end{proof}
